\chapter{State Machine Design and FIFO Queues}
\label{ch:fsmdesign}

In chapter~\ref{ch:sync} we met finite-state machines and analyzed an example
detector. Now we will learn to design them systematically: from a verbal
description to a circuit, and we will see that the same recipe works for the
control part of any device. In the second half of the chapter we will use the
same techniques to build a FIFO queue, which is needed in any system where the
speeds of the data source and the data receiver differ.

The question of the chapter: how do we derive a machine from a verbal
requirement instead of guessing it, so that no case is forgotten? Along the way
we will assemble the main stage of the running project: the mode machine of the
``Security Panel''.

\begin{buildbox}
The \emph{Learn} tab, the ``Circuits and the board'' book, the ``State
machines'' chapter, the lesson ``A state machine with two states'' (then ``The
skeleton of a machine''). Build a machine with the states ``off'' and ``on'',
switched by a press. In the \emph{Waveforms} tab make sure that the state
changes only on the clock edge. Write down on a sheet of paper how you decided
\emph{how many} states are needed---this is the first of the six steps
below.
\end{buildbox}

\section{Six Steps of Design}

\begin{enumerate}
  \item \textbf{Verbal description.} What the device must do, what inputs and
        outputs it has, and when the output must change.
  \item \textbf{State graph.} Draw circles (states) and arrows (transitions)
        labeled ``condition / output''.
  \item \textbf{Transition table.} For every state and every value of the
        inputs, write the next state and the output. No cell may be left
        empty.
  \item \textbf{State encoding.} Assign binary codes to the states
        or use one-hot encoding.
  \item \textbf{Logic.} Obtain the next-state and output functions
        (by hand with Karnaugh maps, or with a program).
  \item \textbf{Verification.} Go through scenarios in simulation, including
        boundary ones: a reset in the middle of operation, extra inputs,
        simultaneous events.
\end{enumerate}

\begin{diybox}[Six Steps on a Small Example]
\textbf{Task.} Go through all six steps for a ``lamp with a button'' machine:
one press turns it on, the next turns it off.\\
\textbf{In Logra.} The \emph{Learn} tab, the ``Circuits and the board'' book,
the ``State machines'' chapter; the \emph{Designer} and \emph{Waveforms}
tabs.\\
\stars{1} \emph{Repeat.} The lesson ``A state machine with two states''.\\
\stars{2} \emph{Change.} The lesson ``The skeleton of a machine'': add a third
state, ``blinking'', which the lamp enters after a double press.\\
\stars{3} \emph{Project.} A pedestrian traffic light (the lesson ``A pedestrian
traffic light''): draw the graph and the transition table \emph{before}
building; when building, no cell of the table may be left empty.\\
\textbf{Expected result.} The waveform repeats the path through the graph: the
states change in the order you drew.\\
\textbf{How to check.} \emph{Waveforms} tab: compare every clock cycle with
the transition table; \emph{RTL} tab: find the state register and the
transition logic.
\end{diybox}

\section{Example: A Vending Machine}

The machine sells an item for 15 units. It accepts coins of 5 and 10 units, one
per clock cycle. When 15 units or more have been inserted, it outputs the
signal \code{vend} \emph{in the same clock cycle}, and then starts counting
again (for simplicity, we give no change).

There are three states: 0, 5, or 10 units inserted. The output depends not only
on the state but also on the coin that has just arrived: a 5 in the ``10
units'' state immediately produces a sale. So this is a \term{Mealy} machine.
Its graph is in Fig.~\ref{fig:vend_graph}: the arrows are labeled with the
inputs and the output, separated by a slash.

\begin{figure}[H]
\centering
\begin{tikzpicture}[scale=0.85, transform shape, font=\sffamily\small, st/.style={circle, draw=FpgaBlue, fill=FpgaLight, minimum size=16mm, thick, align=center}, lbl2/.style={font=\sffamily\scriptsize, fill=white, inner sep=1.5pt}]
  \node[st] (s0) at (0,0) {$S_0$\\{\scriptsize 0 units}};
  \node[st] (s5) at (5,2.6) {$S_5$\\{\scriptsize 5 units}};
  \node[st] (s10) at (10,0) {$S_{10}$\\{\scriptsize 10 units}};
  \draw[-{Stealth}] (s0) -- node[lbl2, above left] {c5/0} (s5);
  \draw[-{Stealth}] (s5) -- node[lbl2, above right] {c5/0} (s10);
  \draw[-{Stealth}] (s5) to[bend right=30] node[lbl2, below left, xshift=-2mm] {c10/1} (s0);
  \draw[-{Stealth}] (s0) to[bend left=18] node[lbl2, above] {c10/0} (s10);
  \draw[-{Stealth}] (s10) to[bend left=18] node[lbl2, below] {c5 or c10/1} (s0);
  \draw[-{Stealth}] (s0) to[loop left] node[lbl2, left] {no coin/0} (s0);
  \draw[-{Stealth}] (s5) to[loop above] node[lbl2, above] {no coin/0} (s5);
  \draw[-{Stealth}] (s10) to[loop right] node[lbl2, right] {no coin/0} (s10);
\end{tikzpicture}
\caption{The graph of a Mealy machine. The label is ``input/output'': after the slash is the value of \code{vend}. All transitions that end in a sale return the machine to $S_0$}
\label{fig:vend_graph}
\end{figure}

The same information is collected in the transition table. It must not have
empty cells:

\begin{table}[H]
\centering
\caption{The transition table of the vending machine ($c5$, $c10$ are coins, \code{vend} is the output)}
\small
\begin{tabular}{lcccc}
\toprule
\textbf{State} & \textbf{No coin} & \textbf{c5} & \textbf{c10} & \textbf{Output} \\
\midrule
$S_0$ (0) & $S_0$ & $S_5$ & $S_{10}$ & 0 \\
$S_5$ (5) & $S_5$ & $S_{10}$ & $S_0$ & 1 only on c10 \\
$S_{10}$ (10) & $S_{10}$ & $S_0$ & $S_0$ & 1 on c5 and on c10 \\
\bottomrule
\end{tabular}
\end{table}

\subsection{How the Machine Works, Cycle by Cycle}

Let us check the machine by hand: feed in the coins $5, 5, 5, 10, 5$ and write
into a table what happens in every clock cycle. The state and the coin
determine both the output and the next state.

\begin{table}[H]
\centering
\caption{A trace of the vending machine}
\small
\begin{tabular}{cllcl}
\toprule
\textbf{Cycle} & \textbf{State} & \textbf{Coin} & \textbf{\code{vend}} & \textbf{Next state} \\
\midrule
0 & $S_0$ & 5 & 0 & $S_5$ \\
1 & $S_5$ & 5 & 0 & $S_{10}$ \\
2 & $S_{10}$ & 5 & \textbf{1} & $S_0$ \\
3 & $S_0$ & 10 & 0 & $S_{10}$ \\
4 & $S_{10}$ & 5 & \textbf{1} & $S_0$ \\
5 & $S_0$ & none & 0 & $S_0$ \\
\bottomrule
\end{tabular}
\end{table}

Note cycles 2 and 4: the signal \code{vend} is one \emph{in the same clock
cycle} in which the last coin was inserted, not in the next one. The output is
computed from the state and the input together, and this is the main hallmark
of a Mealy machine. A Moore machine would dispense the item only in the next
cycle, when the state register has already changed its value.

\subsection{Rules for Reliable Description}

Every machine is divided into three parts: the state register, the
combinational transition logic, and the combinational output logic. Three
techniques protect against typical mistakes.

\begin{itemize}
  \item For every state, a ``next state by default = current state'' is defined
        in advance: if no transition condition fires, the machine stays in
        place, and no unwanted latch appears in the circuit.
  \item For unused state codes, provide a transition ``in all other cases'': if
        the register somehow ends up in such a code (for example, because of a
        fault), the machine returns to the initial state.
  \item Make the reset synchronous: it acts on the clock and therefore causes
        no timing problems.
\end{itemize}

How these three rules look in code, and how this machine is checked on a
sequence of coins, is shown in Section~\ref{sec:v-fsm3}.

Let us compare the two kinds of machines.

\begin{table}[H]
\centering
\caption{Moore and Mealy machines}
\small
\begin{tabularx}{\textwidth}{l X X}
\toprule
 & \textbf{Moore} & \textbf{Mealy} \\
\midrule
The output depends on & the state & the state and the inputs \\
Number of states & usually larger & usually smaller \\
Reaction to an input & one cycle later & immediately, in the same cycle \\
Noise on the input & does not reach the output & reaches the output \\
\bottomrule
\end{tabularx}
\end{table}

\begin{diybox}[The Vending Machine]
\textbf{Task.} Check the vending machine on coin sequences and teach it to
handle simultaneous coins.\\
\textbf{In Logra.} The \emph{Examples} tab, project \code{vending-fsm}
(the machine counts events and outputs a result after a set number);
the \emph{Waveforms} and \emph{RTL} tabs.\\
\stars{1} \emph{Repeat.} Compare one long button hold with three separate
presses and releases: on the waveform find when the dispense pulse
appears.\\
\stars{2} \emph{Change.} In the \emph{Code} tab increase the number of events
before dispensing from three to four and check the dispense pulse on the
waveform. Then use the table of this chapter to make a trace for the coins
$10, 5, 10$.\\
\stars{3} \emph{Project.} Problem~2 of this chapter: add a column ``both coins
at once'' to the transition table and check the table for empty cells.\\
\textbf{Expected result.} The state accumulates events, and after the set
number it outputs a one-cycle pulse and returns to the start.\\
\textbf{How to check.} \emph{Waveforms} tab: between dispense pulses there are
exactly as many events as you set; \emph{RTL} tab: find the register of the
accumulated sum.
\end{diybox}

\section{A Three-Block Machine in Verilog}
\label{sec:v-fsm3}

A state graph and a transition table translate into Verilog almost
mechanically. A machine is written as three separate parts: each answers its
own question, and they must not be mixed, or latches and hard-to-find errors
appear in the circuit.

\begin{table}[H]
\centering
\caption{The three blocks of a machine and how they are written}
\small
\begin{tabularx}{\textwidth}{l X X}
\toprule
\textbf{Block} & \textbf{What it does} & \textbf{How it is written} \\
\midrule
1. State register & remembers where we are; reset sets the initial state & \code{always @(posedge clk)}, assignment \code{<=}, reset first \\
2. Transition logic & computes the next state \code{next} from the state and inputs & \code{always @(*)}, \code{next = state;} at the start, then \code{case (state)} with a \code{default} branch \\
3. Output logic & Moore: the output depends only on \code{state}. Mealy: on \code{state} and inputs & \code{always @(*)} with \code{case}, or \code{assign}; assignment \code{=} \\
\bottomrule
\end{tabularx}
\end{table}

\subsection{A Traffic Light: A Moore Machine}

Three states, in each one lamp is lit; the time in a state is counted by the
clock counter \code{t}. The durations are parameters, so green can be changed
from 4 clocks to 6 with a single number.

\vfile{\codepath/b9a_fsm/traffic3.v}{code/b9a\_fsm/traffic3.v}

\textbf{How to read it.} Block~1 is the only place where \code{state} and
\code{t} change: on the clock edge the state takes the value \code{next}. In
block~2 the first line \code{next = state} sets ``stay'' as the default, and the
signal \code{last} (``this is the last clock of the state'') decides whether
to leave. The \code{default} branch returns the machine to the start if the
register falls into the unused code \code{3}. Block~3 depends only on
\code{state}: this is exactly a Moore machine, so the lamp changes exactly
together with the state.

\shot{br-fsm-wave}{The traffic light on a waveform: one lamp is lit in each state ($0x1$ is green, $0x2$ is yellow, $0x4$ is red); a screenshot of the \code{traffic\_fsm} module from \code{code/00\_theory}, \code{traffic3} gives the same picture}{fig:b-fsm-wave}

The testbench \code{b9a\_fsm/tb\_traffic3.v} checks three rounds by formula:
green 4 clocks, yellow 2, red 4. It prints \code{traffic3: OK}.

\begin{diybox}[A Traffic Light by Parameters]
\textbf{Task.} Run the traffic light and change the durations.\\
\textbf{In Logra.} The Code tab: paste \code{traffic3.v}, ``Synthesize''
(Ctrl+Enter), the RTL and Waveforms tabs (30 cycles, a \code{rst} pulse at the
start).\\
\textbf{Expected result.} \code{lamp} takes the values \code{001}, \code{010},
\code{100} for 4, 2 and 4 clocks and repeats.\\
\textbf{How to check.} On RTL find the register \code{state}, the register
\code{t} and the output multiplexer. In a terminal:
\code{iverilog -g2012 -o tb tb\_traffic3.v traffic3.v}, then \code{vvp tb}.
\stars{2} Add a fourth state: before green, red and yellow are lit together for
1 clock (\code{3'b110}).\\
\textbf{Sample answer.} You need a new constant, a branch in blocks~2 and~3, and
a changed transition out of red.
\end{diybox}

\begin{vcode}
module traffic4 (
    input            clk, rst,
    output reg [2:0] lamp
);
    localparam [1:0] RED = 2'd0, RED_YEL = 2'd1, GREEN = 2'd2, YELLOW = 2'd3;
    reg [1:0] state, next;
    reg [2:0] t;
    wire last = (state == RED && t == 3'd3) || (state == RED_YEL && t == 3'd0)
             || (state == GREEN && t == 3'd3) || (state == YELLOW && t == 3'd1);
    always @(posedge clk) begin
        if (rst) begin state <= RED; t <= 3'd0; end
        else begin state <= next; t <= last ? 3'd0 : t + 3'd1; end
    end
    always @(*) begin
        next = state;
        if (last) case (state)
            RED:     next = RED_YEL;
            RED_YEL: next = GREEN;
            GREEN:   next = YELLOW;
            default: next = RED;
        endcase
    end
    always @(*) begin
        case (state)
            RED:     lamp = 3'b100;
            RED_YEL: lamp = 3'b110;
            GREEN:   lamp = 3'b001;
            default: lamp = 3'b010;
        endcase
    end
endmodule
\end{vcode}


\subsection{The Vending Machine: A Mealy Machine}

The machine from the beginning of the chapter is written with the same three
blocks. The only difference is in block~3: the output \code{vend} is computed
from the state \emph{and from the inputs} \code{c5}, \code{c10}.

\vfile{\codepath/00_theory/vending.v}{code/00\_theory/vending.v}

\textbf{How to read it.} The states are how many rubles have been inserted. In
block~2 both coins are described for every state, and when there is no coin
\code{next = state} stays. Block~3 is a single formula: a sale if a ten arrives
in $S_5$ or any coin arrives in $S_{10}$. All three rules from ``Rules for
Reliable Description'' are in place: a default value, \code{default}, and a
synchronous reset.

\shot{n-vend-wave}{A Mealy machine: \code{vend} rises in the same clock in which the last coin arrives}{fig:n-vend-wave}

On the Waveforms tab the coins are set in the ``Events'' section. The signal
\code{vend} rises inside the clock, as soon as the ``deciding'' coin appears at
the input, without waiting for the next edge---this is what makes it a Mealy
machine. If a ``clean'' output is needed, it is passed through a flip-flop and
lags by one clock. The testbench \code{00\_theory/tb\_vending.v} throws the same
coins as the trace in the table of the chapter (5+5+5, 10+5, 5+10, 10+10) and
counts the sales: there must be four.

\begin{diybox}[The Vending Machine in Code]
\textbf{Task.} Run the machine and teach it a new rule.\\
\textbf{In Logra.} The Code tab: \code{vending.v}, ``Synthesize'', Waveforms:
the ``Events'' section, pulses \code{c5} and \code{c10} by clock.\\
\textbf{Expected result.} For the coins 5, 5, 5, \code{vend = 1} in the clock of
the third coin; for the coins 10, 5, in the clock of the five.\\
\textbf{How to check.} Run \code{tb\_vending.v}: \code{vending: OK}.
\stars{2} Make the price 20 rubles (add a state $S_{15}$).
\end{diybox}

\textbf{Sample answer.}
\begin{vcode}
module vending20 (
    input  clk, rst, c5, c10,
    output vend
);
    localparam [1:0] S0 = 2'd0, S5 = 2'd1, S10 = 2'd2, S15 = 2'd3;
    reg [1:0] state, next;
    always @(posedge clk) begin
        if (rst) state <= S0;
        else     state <= next;
    end
    always @(*) begin
        next = state;
        case (state)
            S0:  if (c10) next = S10; else if (c5) next = S5;
            S5:  if (c10) next = S15; else if (c5) next = S10;
            S10: if (c10) next = S0;  else if (c5) next = S15;
            S15: if (c10 | c5) next = S0;
        endcase
    end
    assign vend = (state == S10 && c10) || (state == S15 && (c5 | c10));
endmodule
\end{vcode}


\subsection{The 1-0-1 Sequence Detector}

The detector from Chapter~\ref{ch:sync} and from the panel code is also a Mealy
machine. The states: \code{S0}---``nothing useful'', \code{S1}---``saw 1'',
\code{S10}---``saw 10''. The output \code{hit} is one when another one arrives
in state \code{S10}. A one in any state can start a new sequence, so the
triples overlap: the stream \code{10101} gives two hits.

\vfile{\codepath/b9a_fsm/seq101.v}{code/b9a\_fsm/seq101.v}

\textbf{How to read it.} Block~2 holds all the ``memory'' of the sequence: from
\code{S10} on input 1 the machine goes not to \code{S0} but to \code{S1}, because
the one that has just arrived has already started a new triple. Block~3 is one
line, a Mealy output.

The testbench \code{tb\_seq101.v} feeds the stream \code{1 0 1 0 1 1 0 1}, which
has three hits, and the stream \code{1 1 0 0 0 1 1 1}, which has none. The
bit-feeding task \code{send(b)} looks at \code{hit} \emph{before} the clock edge:
this is how a Mealy output is checked.

\begin{diybox}[Your Own Detector]
\textbf{Task.} Run the detector and rebuild it for a different code.\\
\textbf{In Logra.} The Code tab: \code{seq101.v}; Waveforms: the ``Events''
section, the signal \code{x} by clock \code{1, 0, 1, 0, 1}. Or the lesson ``Code
lock'' (the Learn tab, the Verilog book).\\
\textbf{Expected result.} Two \code{hit} pulses on the stream \code{10101}; none
on \code{110} and \code{011}.\\
\textbf{How to check.} \code{tb\_seq101.v} prints \code{seq101: OK}.
\stars{2} Rebuild the detector for the code \code{1-1-0}.\\
\textbf{Sample answer.} The transitions and the output change: ``saw 1'', ``saw
11'', and a hit on a zero after two ones.
\end{diybox}

\begin{vcode}
module seq110 (
    input  clk, rst, x,
    output hit
);
    localparam [1:0] S0 = 2'd0, S1 = 2'd1, S11 = 2'd2;
    reg [1:0] state, next;
    always @(posedge clk) begin
        if (rst) state <= S0;
        else     state <= next;
    end
    always @(*) begin
        next = state;
        case (state)
            S0:      next = x ? S1 : S0;
            S1:      next = x ? S11 : S0;
            S11:     next = x ? S11 : S0;   // 1-1-1 still contains 1-1
            default: next = S0;
        endcase
    end
    assign hit = (state == S11) && !x;
endmodule
\end{vcode}


\section{State Encoding and Reduction}

Three states need two flip-flops with binary encoding
(\code{00}, \code{01}, \code{10}) and three with one-hot encoding. The choice
of encoding can change the size and speed of the circuit. The
synthesizer (Chapter~\ref{ch:flow}) can usually re-encode a machine by itself
if you describe it using a template with named states; still, it is useful to
be able to choose the encoding by hand.

Sometimes two states behave identically: for any inputs they give the same
outputs and move to the same states. Such states can be merged into one, and
the machine becomes simpler. The test is simple: if all rows of the transition
table of two states coincide (up to the names of these same states), they can
be merged.

Let us show three encodings on one machine. Four states \code{A}, \code{B},
\code{C}, \code{D} go round in a circle, the transition is made on the signal
\code{go}, and the output \code{y} is one in states \code{C} and \code{D}. In
binary code they are written as the numbers 0--3, in Gray code neighbouring
states differ in one bit (\code{00, 01, 11, 10}), in one-hot code every state
has its own flip-flop and exactly one is set.

\vfile{\codepath/b9a_fsm/state_enc.v}{code/b9a\_fsm/state\_enc.v}

\textbf{How to read it.} The first two modules have an identical \code{case};
only the constants \code{A\ldots D} differ: the choice of code changes nothing in
the text of the machine. But in Gray code the output turned out to be simply the
high bit \code{st[1]}, and the output logic disappeared. The one-hot module has
no \code{case}: a transition is a shift of the one to the neighbouring
flip-flop, and the output is an OR of two flip-flops. The price is 4
flip-flops instead of 2; the gain is almost no logic between flip-flops. The
testbench \code{tb\_state\_enc.v} feeds the same \code{go} to all three modules
and checks that the outputs match on every clock.

\begin{diybox}[Binary and One-Hot Encoding]
\textbf{Task.} Build the same machine in two ways and compare the cost of
each.\\
\textbf{In Logra.} The \emph{Designer}, \emph{Waveforms}, \emph{RTL}, and
\emph{Gowin} tabs.\\
\stars{1} \emph{Repeat.} For a ``red---green---yellow'' machine write out the
binary state codes and the one-hot ones.\\
\stars{2} \emph{Change.} Build the machine in the Designer on two flip-flops
(binary encoding), then on three (one-hot); the output waveforms must
match.\\
\stars{3} \emph{Project.} Compare the number of flip-flops and logic gates in
the \emph{RTL} and \emph{Gowin} tabs; write your conclusion in one line.\\
\textbf{Expected result.} The one-hot version uses more flip-flops and less
logic per transition.\\
\textbf{How to check.} \emph{Waveforms} tab: the outputs of both versions are
identical in all clock cycles; \emph{RTL} tab: count the flip-flops.\\
\textbf{The same in Verilog.} The \emph{Code} tab: paste \code{state\_enc.v},
``Synthesize'', make \code{ring\_bin}, \code{ring\_gray}, \code{ring\_onehot} the
top module in turn and compare the number of registers on RTL (2, 2 and 4).
\stars{2} Complete a one-hot ring for the three traffic-light states.
\end{diybox}

\textbf{Sample answer.}
\begin{vcode}
module ring3_onehot (
    input      clk, rst, go,
    output     green, yellow, red
);
    reg [2:0] st;                       // st[0] = green, st[1] = yellow, st[2] = red
    always @(posedge clk) begin
        if (rst)     st <= 3'b001;
        else if (go) st <= {st[1:0], st[2]};
    end
    assign green  = st[0];
    assign yellow = st[1];
    assign red    = st[2];
endmodule
\end{vcode}


\section{Running Project: The ``Security Panel'' Machine}
\label{sec:panel-fsm}

So far the project has had separate parts: a sensor, an ``alarm'' logic, alarm
memory, and a time counter. Now we will assemble them into a device with modes.
The panel is in one of three states: \term{DISARMED} (the sensors are ignored,
the residents are home), \term{ARMED} (the door sensor is watching), and
\term{ALARM} (the siren is on). The modes are switched by entering a code: bits
are fed to the button on successive clock cycles, and the code counts as entered
if \code{1-0-1} arrives in a row.

Here two machines are connected by one signal. The first is the code detector
from chapter~\ref{ch:sync}: it watches the input and outputs $c=1$ in the
clock cycle when the code has been entered correctly. The second is the mode
machine, which reads the signal $c$ and the door sensor $d$ (1 means the door
is open).

\begin{figure}[H]
\centering
\begin{tikzpicture}[
  st/.style={circle, draw=FpgaDark, very thick, minimum size=19mm, font=\sffamily\small, align=center, fill=FpgaLight},
  tr/.style={-{Stealth[length=3mm]}, thick, FpgaDark},
  lb/.style={font=\sffamily\scriptsize, fill=white, inner sep=1.5pt}]
  \node[st] (off) at (0,0) {DISARMED\\\scriptsize siren 0};
  \node[st] (arm) at (5.2,0) {ARMED\\\scriptsize siren 0};
  \node[st, fill=red!12] (al) at (10.4,0) {ALARM\\\scriptsize siren 1};
  \draw[tr] (off) to[bend left=22] node[lb, above] {$c=1$} (arm);
  \draw[tr] (arm) to[bend left=22] node[lb, below] {$c=1$} (off);
  \draw[tr] (arm) -- node[lb, above] {$d=1$, $c=0$} (al);
  \draw[tr] (al) to[bend left=50] node[lb, below] {$c=1$} (off);
  \draw[tr] (off) edge[loop above, looseness=7] node[lb, above] {$c=0$} (off);
  \draw[tr] (arm) edge[loop above, looseness=7] node[lb, above] {$d=0$, $c=0$} (arm);
  \draw[tr] (al) edge[loop above, looseness=7] node[lb, above] {$c=0$} (al);
  \draw[tr, FpgaRed] (-1.9,-1.9) node[left, font=\sffamily\scriptsize, text=FpgaRed] {reset} -- (off.south west);
\end{tikzpicture}
\caption{The panel mode machine (Moore). The siren is on only in the ALARM
state. The condition for leaving any mode for DISARMED is the correct code}
\label{fig:panel_fsm}
\end{figure}

Let us check against the six steps. The states are given (steps 1--2). The
transition table (step 3) has no empty cells: for every state, all four
combinations $(d,c)$ are described. In the ARMED state with $d=1$ and $c=1$ the
code wins: the owner came in and entered the code, so no siren is needed.

\begin{table}[H]
\centering
\caption{The transition table of the panel machine ($d$ is door open, $c$ is correct code)}
\small
\begin{tabular}{l|cccc|c}
\toprule
\textbf{State} & $d{=}0,c{=}0$ & $d{=}0,c{=}1$ & $d{=}1,c{=}0$ & $d{=}1,c{=}1$ & \textbf{Siren} \\
\midrule
DISARMED & DISARMED & ARMED & DISARMED & ARMED & 0 \\
ARMED    & ARMED & DISARMED & ALARM & DISARMED & 0 \\
ALARM    & ALARM & DISARMED & ALARM & DISARMED & 1 \\
\bottomrule
\end{tabular}
\end{table}

Now let us write both machines in Verilog. The code detector is the same
\code{seq101} (in the file it is renamed \code{code\_detector} so that the file
is self-contained). The mode machine has three blocks: the register
\code{mode}, the transitions from the table above, and the Moore output
\code{siren}. The top module \code{panel\_fsm} only connects them with the
signal \code{c}.

\vfile{\codepath/b9a_fsm/panel_fsm.v}{code/b9a\_fsm/panel\_fsm.v}

\textbf{How to read it.} In the transition block the \code{ARMED} branch checks
\code{c} first and only then \code{d}: this way in the cell $d{=}1,c{=}1$ the
code wins, as the table requires. The \code{ALARM} state is left only by the
code: the alarm does not stop by itself, even if the door is closed. The
testbench \code{tb\_panel\_fsm.v} goes through the four items of the acceptance
criterion from ``Do it yourself'' below and prints \code{panel\_fsm: OK}.

\begin{diybox}[Project Stage 9: The Mode Machine]
\textbf{Task.} Build the ``Security Panel'' machine with entry of the code
\code{1-0-1} and a siren.\\
\textbf{In Logra.} The \emph{Learn} tab, the ``Circuits and the board'' book,
the ``State machines'' chapter: the lesson ``A code lock'' (code
\code{1, 1, 0}); the \emph{Code} tab, the ``Verilog and the board'' book: the
lesson ``A code lock: 1, 0, 1''; the \emph{Examples} tab: the project
\code{door-controller} as a model of a machine with modes; the
\emph{Designer}, \emph{Waveforms}, and \emph{RTL} tabs.\\
\stars{1} \emph{Repeat.} The ``A code lock'' lesson from blocks: build the
lock and make sure it opens only with the right triple.\\
\stars{2} \emph{Change.} Rebuild the detector for the code \code{1-0-1}, then
add the input $d$ and the output ``siren'' to the mode machine according to the
table of this chapter.\\
\stars{3} \emph{Project.} \textbf{A stage of the ``Security Panel'' project:}
a machine DISARMED $\to$ ARMED $\to$ ALARM with entry of the code \code{1-0-1}
and alarm reset by the same code. Acceptance criterion: on the waveform (a)
after \code{101} is entered the panel moves to ARMED, (b) opening the door
($d=1$) in the ARMED mode turns the siren on in the next clock cycle, (c) in
the DISARMED mode opening the door does \emph{not} turn the siren on, (d)
entering \code{101} in the ALARM state turns the siren off.\\
\textbf{Expected result.} All four points of the criterion hold on the
waveform, and the transition table has no empty cells.\\
\textbf{How to check.} \emph{Waveforms} tab: in turn, feed the code,
open the door, and enter the code again; \emph{RTL} tab: find the state
register and the two logic blocks (transitions and output).\\
\textbf{The same in Verilog.} The \emph{Code} tab: paste \code{panel\_fsm.v},
``Synthesize'', on Waveforms feed \code{code\_bit} on the clocks \code{1, 0, 1},
then \code{d = 1}. The sample answer to \stars{3} is the file
\code{panel\_fsm.v} itself; the testbench \code{tb\_panel\_fsm.v} checks
items (a)--(d).
\end{diybox}

\section{The FIFO Queue}

Often the source of data and its consumer work at different speeds: a sensor
delivers samples in bursts, while a communication line sends them slowly. We
need a buffer that stores the data and outputs it in the same order. Such a
buffer is called a \term{FIFO} (first in---first out). It works like a queue
in a store.

\subsection{The Ring Buffer}

Inside a FIFO there is ordinary memory (an array of words) and two pointers:
the \emph{write pointer} (where to put the next word) and the \emph{read
pointer} (where to take from). Each operation moves its pointer by one, and at
the end of the memory it returns to zero, like the hand of a clock. The number
of words in the queue is the number accumulated between the pointers: when it
is zero, the queue is empty (\code{empty}); when it equals the size of the
memory, it is full (\code{full}).

\begin{figure}[H]
\centering
\begin{tikzpicture}[font=\sffamily\small]
  \foreach \i in {0,...,7} {
    \pgfmathsetmacro{\a}{90-\i*45}
    \node[draw, minimum size=9mm, fill=white] (c\i) at (\a:2.2) {\i};
  }
  % occupied cells: 2,3,4
  \foreach \i in {2,3,4} \node[draw, minimum size=9mm, fill=FpgaLight] at (90-\i*45:2.2) {\i};
  \node[font=\sffamily\bfseries, text=FpgaGreen!50!black] at (90-2*45:3.5) {rp};
  \draw[-{Stealth}, FpgaGreen!50!black, thick] (90-2*45:3.3) -- (90-2*45:2.8);
  \node[font=\sffamily\bfseries, text=FpgaOrange] at (90-5*45:3.95) {wp};
  \draw[-{Stealth}, FpgaOrange, thick] (90-5*45:3.55) -- (90-5*45:2.85);
  \node[align=left, anchor=west] at (4.4,1.5) {3 words are stored (cells 2, 3, 4).\\Reading will start from cell 2.\\The next write goes to cell 5.};
  \node[align=left, anchor=west] at (4.4,-1.2) {After cell 7 the pointers\\return to 0: that is why for\\8 words 3-bit counters\\are enough.};
\end{tikzpicture}
\caption{A ring buffer: the cells between the read pointer \code{rp} and the write pointer \code{wp} are filled}
\end{figure}

\subsection{The full and empty Flags}

For a queue of 8 words, 3-bit pointers are enough, because a 3-bit counter
``wraps around'' by itself after the value 7. A separate 4-bit word counter
\code{count} (values from 0 to 8) simplifies the flags \code{empty} and
\code{full}. Writing is allowed only while the queue is not full, and reading
only while it is not empty.

\begin{table}[H]
\centering
\caption{A trace of the queue: three words written, then three read}
\small
\begin{tabular}{lcccc}
\toprule
\textbf{Action} & \textbf{wp} & \textbf{rp} & \textbf{count} & \textbf{empty} \\
\midrule
after reset & 0 & 0 & 0 & 1 \\
wrote 11 & 1 & 0 & 1 & 0 \\
wrote 22 & 2 & 0 & 2 & 0 \\
wrote 33 & 3 & 0 & 3 & 0 \\
read (got 11) & 3 & 1 & 2 & 0 \\
read (got 22) & 3 & 2 & 1 & 0 \\
read (got 33) & 3 & 3 & 0 & 1 \\
\bottomrule
\end{tabular}
\end{table}

The words come out in the order in which they went in. Note that the equality
of the pointers $\text{wp}=\text{rp}$ does not distinguish an empty queue from
a full one: in both cases the pointers coincide. So the flags are taken not from
a comparison of the pointers but from the word counter:
$\text{empty}=(\text{count}=0)$, $\text{full}=(\text{count}=8)$.

\subsection{Overflow and Underflow on Waveforms}

What happens if the source writes while the queue is already full, or the
receiver reads from an empty queue? Consider a queue of 4 words. For five
clock cycles in a row we try to write, then for five cycles in a row we read.

\begin{figure}[H]
\centering
\begin{tikztimingtable}[timing/slope=0.05, xscale=1.1, yscale=1.4, timing/font=\sffamily\small, timing/rowdist=3.2ex]
  CLK    & 10{1L 1H} \\
  wr     & 10H 10L \\
  rd     & 10L 10H \\
  count  & 2D{0} 2D{1} 2D{2} 2D{3} 2D{4} 2D{4} 2D{3} 2D{2} 2D{1} 2D{0} \\
  full   & 8L 4H 8L \\
  empty  & 2H 16L 2H \\
\end{tikztimingtable}
\caption{A queue of 4 words. In the fifth cycle the queue is already full (\code{full}=1): the fifth write
is ignored and the word is lost. In the last cycle the queue is empty (\code{empty}=1): the tenth
read outputs nothing, and the counter stays at zero}
\label{fig:fifo_flags}
\end{figure}

The queue ``protects'' itself: when $\text{full}=1$ a write does not move the
pointer, and when $\text{empty}=1$ a read does not move its pointer. But a word
that arrives at a full queue is \emph{lost}: only the writer can fix this, by
watching \code{full} and waiting. So in a system the source must either check
the flag or have a margin in the size of the queue.

\subsection{Example: A Four-Word Queue}

\begin{buildbox}[Watch the Queue at Work]
The \emph{Examples} tab, the project \code{fifo-four} (``Four-word FIFO'').
After a reset the internal source writes the words 1, 2, 3, 4 by itself and
stops when the queue is full. Each button press reads the oldest word. The
pointers \code{write\_ptr} and \code{read\_ptr} move independently, and the
counter \code{used} limits access. On the \emph{Waveforms} tab follow how the
pointers ``wrap around''.
\end{buildbox}

\subsection{The Queue in Verilog: A Word Counter}

The same buffer for eight words of eight bits. The array \code{mem} is the
memory, \code{wp} and \code{rp} are the write and read pointers, and
\code{count} is the word counter from which the flags are taken.

\vfile{\codepath/00_theory/fifo8.v}{code/00\_theory/fifo8.v}

\textbf{How to read it.} The conditions \code{do\_wr} and \code{do\_rd} switch on
the protection: a write to a full queue and a read from an empty one change
nothing. The three-bit \code{wp} and \code{rp} ``wrap around'' by themselves
after seven, so no separate check for the end of memory is needed. The line
\code{count <= count + do\_wr - do\_rd} handles a simultaneous write and read:
the counter does not change. The synthesizer puts the array \code{mem} in
the chip's memory or in cells.

\shot{n-fifo-wave}{The queue at work: the words \code{0B}, \code{16}, \code{21} are written, then read in the same order}{fig:n-fifo-wave}

On the waveform \code{count} grows with each write and falls with each read,
\code{empty} is cleared after the first write, and the words come out on
\code{dout} in the order of writing. The testbench \code{00\_theory/tb\_fifo8.v}
writes ten words (the last two do not fit), checks that \code{full = 1}, reads
eight words in order, and makes sure the queue is empty.

\subsection{The Queue in Verilog: An Extra Pointer Bit}

The word counter can be left out. If the pointers are made \emph{one bit wider}
than the address, the high bit counts ``laps'': the queue is empty when the
pointers are equal as a whole, and full when the low bits (the address) are
equal but the high bits differ (the writer has gone exactly one lap ahead of
the reader). This is how queues with parameters are written: depth
$2^{\text{AW}}$, word width \code{DW}.

\vfile{\codepath/b9a_fsm/fifo_ptr.v}{code/b9a\_fsm/fifo\_ptr.v}

\textbf{How to read it.} The memory address is taken from \code{wp[AW-1:0]} and
\code{rp[AW-1:0]}, and the flags from the full pointers. The word \code{dout}
is shown at once (without waiting for a clock): it is always equal to the oldest
word while the queue is not empty. The testbench \code{tb\_fifo\_ptr.v} compares
the queue with a simple model (an array and two integer indices) on 12 writes in
a row, 12 reads, and 400 random clocks: \code{full}, \code{empty} and the order
of words are all checked.

\begin{diybox}[The Queue in Code]
\textbf{Task.} Run the queue and catch an overflow.\\
\textbf{In Logra.} The Code tab: \code{fifo\_ptr.v} (or \code{fifo8.v}),
``Synthesize'', Waveforms: the ``Events'' section, \code{wr} pulses by clock
with \code{din}, then \code{rd}.\\
\textbf{Expected result.} After eight writes \code{full = 1}, the ninth write
changes nothing; the words come out in the order of writing; after the last
read \code{empty = 1}.\\
\textbf{How to check.} \code{tb\_fifo\_ptr.v} prints \code{fifo\_ptr: OK}.
\stars{2} Make a queue of 4 words of 4 bits without changing the module text.\\
\textbf{Sample answer.} The parameters are changed when the instance is created:
\code{fifo\_ptr \#(.DW(4), .AW(2)) u (.clk(clk), .rst(rst), .wr(wr), .rd(rd),
.din(din), .dout(dout), .empty(empty), .full(full));}. The module text itself
stays the same; with \code{.AW(2)} the flag \code{full} comes after four writes
(to check this with the testbench, replace the number \code{8} with \code{4} in
\code{tb\_fifo\_ptr.v} in the check and in the model).
\end{diybox}

\lookahead{in large projects a queue can be huge and keep its words in the block
memory of the chip; there is also a \emph{dual-clock} queue, where writing runs on
one clock signal and reading on another---then the pointers are passed between
domains in Gray code (chapter~\ref{ch:clocks}).}

\begin{warnbox}[Common mistakes]
\begin{itemize}
  \item \textbf{Telling empty from full by pointer equality.} When
        $\text{wp}=\text{rp}$ it is unclear whether there are zero words or
        all of them. You need a word counter or an extra pointer bit.
  \item \textbf{Writing when \code{full} and reading when \code{empty}.} A word
        is lost, or garbage is read. The source and the receiver must look at
        the flags.
  \item \textbf{The pointer does not return to zero.} If a queue of 8 words
        uses a 4-bit pointer without wrapping at the boundary, the address
        will go outside the memory.
  \item \textbf{Forgetting the reset.} After power-up the pointers and the
        counter must be zeros, otherwise the queue ``sees'' words that do not
        exist.
\end{itemize}
\end{warnbox}

\begin{diybox}[A Queue with full and empty Flags]
\textbf{Task.} Follow the operation of a ring buffer and find the moments of
``full'' and ``empty''.\\
\textbf{In Logra.} The \emph{Examples} tab, the project \code{fifo-four};
the \emph{Waveforms} and \emph{RTL} tabs.\\
\stars{1} \emph{Repeat.} Wait until the queue fills with the words
1, 2, 3, 4, then press the button several times; compare the pointers with the
trace table of this chapter.\\
\stars{2} \emph{Change.} Find on the waveform the clock cycle in which \code{used}
reaches 4 and writing stops (this is how the ``full'' flag shows itself), and the cycle
in which \code{used} becomes 0 (``empty''). Problem~5: what would happen without
the protection?\\
\stars{3} \emph{Project.} In the \emph{Code} tab enlarge the queue to eight
words (3-bit pointers, a 4-bit \code{used}, comparison with 8) and
check the full and empty flags.\\
\textbf{Expected result.} When the queue is full, a write changes neither the
pointer nor the counter; the words come out in the order they were written.\\
\textbf{How to check.} \emph{Waveforms} tab: the sequence of
\code{used} during filling and reading matches Fig.~\ref{fig:fifo_flags};
\emph{RTL} tab: find the word memory and the two pointers.
\end{diybox}

\section{Questions and Problems}
\label{sec:tasks-fsmdesign}

The answers are collected in Appendix~\ref{app:answers}.

\begin{taskbox}[Problems for the Chapter]
\begin{enumerate}
  \item \stars{2}~Draw the graph of a Moore machine with three states for a traffic light in which
        red is on for 4 clock cycles, yellow for 1, and green for 3. How many flip-flops
        are needed with binary encoding? With one-hot encoding?
  \item \stars{2}~The vending machine receives the coins $c5$ and $c10$ at the same time (a
        sensor fault). How will the description from this chapter behave? How should the
        transition table be changed to take this into account?
  \item \stars{2}~Draw the transition table of a detector of the sequence \code{110}
        (with overlap). What type of machine results, Moore or Mealy?
  \item \stars{1}~A FIFO queue has 16 words. What width do the pointers and
        the word counter need?
  \item \stars{1}~11 words were written in a row into a queue of 8 words. How many words will be
        lost, and why?
  \item \stars{1}~In what way is a Mealy machine more convenient than a Moore machine, and in what way more dangerous?
  \item \stars{2}~\emph{Find the error.} The designer of a queue of 8 words defined
        $\text{full}=(\text{wp}=\text{rp})$ and $\text{empty}=(\text{wp}=\text{rp})$. In which
        two situations will the flags give a wrong result, and how can they be fixed?
  \item \stars{2}~\emph{Find the error.} In the transition table of the panel machine
        (section ``Running Project'') the designer left the cell
        $d{=}1,c{=}1$ empty for the ARMED state, deciding that ``this never happens''. Describe a
        scenario in which it happens, and say what the empty
        cell will lead to. What transition would you write?
  \item \stars{3}~\emph{Mini-project.} Build the ``Security Panel'' mode machine
        with entry of the code \code{1-0-1} (stage~9). Acceptance criterion: on the waveform,
        after \code{101} the panel is in the ARMED state; $d=1$ in the ARMED mode
        turns the siren on from the next clock cycle; $d=1$ in the DISARMED mode does not
        turn it on; entering \code{101} again in the ALARM state turns it off.
  \item \stars{3}~\emph{Mini-project.} Extend the 4-word queue with the outputs
        \code{full} and \code{empty} and make writing to a full
        queue impossible. Acceptance criterion: with five writes in a row,
        on the waveform \code{used}$=4$, the write pointer has stopped, and the fifth word did not
        overwrite the first; with five reads in a row the fifth read changes nothing.
\end{enumerate}
\end{taskbox}

\section{What's Next}

We have learned to describe the behavior of a device with a state machine and to
store data in a queue, that is, to design a circuit of any complexity and check
it in Logra. But what is such a circuit made of inside the chip? An FPGA has no
separate gates: combinational logic is replaced by lookup tables, a state
machine and a queue are built from flip-flops and memory blocks, and the wires
between them are switched when the configuration is loaded. That is the subject
of Chapter~\ref{ch:inside}: we look inside the chip, and then go through the
board and the flow from text to a bitstream.
